\section{Dokumentation KI-basierter Hilfsmittel}
\label{app:ki_hilfsmittel}

Im Rahmen der prototypischen Realisierung des in dieser Arbeit untersuchten
Transformationssystems wurde ChatGPT von OpenAI als KI-basiertes Hilfsmittel
bei der Softwareentwicklung des SysMLv2-Generators eingesetzt. Der Einsatz
betraf insbesondere die Erstellung und Überarbeitung von Programmcode, die
Ableitung und Ergänzung automatisierter Tests sowie die Analyse technischer
Zwischenstände.

Tabelle~\ref{tab:ki_hilfsmittel} fasst die wesentlichen Einsatzbereiche und
die jeweilige menschliche Kontrolle zusammen.

\begin{table}[htbp]
    \centering
    \caption[Einsatz KI-basierter Hilfsmittel bei der Prototypentwicklung]{
        Einsatz von ChatGPT bei der prototypischen Realisierung des
        SysMLv2-Generators und die jeweils zugehörige menschliche Prüfung.}
    \label{tab:ki_hilfsmittel}
    \begin{tabularx}{\textwidth}{p{3.2cm}p{4.2cm}X}
        \hline
        \textbf{Hilfsmittel} &
        \textbf{Einsatzbereich} &
        \textbf{Verwendung und Kontrolle} \\
        \hline

        ChatGPT, OpenAI &
        Programmcode &
        Unterstützung bei der Erstellung und Überarbeitung einzelner
        Implementierungsbestandteile. Vorgeschlagene Änderungen wurden vor
        ihrer Übernahme gegen die fachliche Spezifikation und die bestehende
        Architektur geprüft. \\

        ChatGPT, OpenAI &
        Automatisierte Tests &
        Unterstützung bei der Ableitung und Ergänzung von Tests für
        spezifizierte technische Eigenschaften und Verarbeitungsschritte.
        Die Akzeptanz der Tests und ihrer erwarteten Ergebnisse erfolgte
        durch den Autor. \\

        ChatGPT, OpenAI &
        Technische Analyse &
        Unterstützung bei der Analyse von Implementierungszuständen,
        Fehlerbildern und möglichen Anpassungen. Architekturentscheidungen
        sowie die Auswahl und Freigabe umgesetzter Änderungen verblieben
        beim Autor. \\

        \hline
    \end{tabularx}
\end{table}

Die KI-generierten Vorschläge besaßen innerhalb des Entwicklungsprozesses
keine eigene Entscheidungsautorität. Fachliche Spezifikation,
Architekturentscheidungen, Bewertung der vorgeschlagenen Änderungen und
deren endgültige Übernahme erfolgten durch den Autor. Die dabei beobachtete
Rolle menschlicher Prüfung wird in
Abschnitt~\ref{sec:reflection_ai_development} reflektiert.